% Source ownership:
% - Itxaso

% ============================================================
\section{Input/Output Module}
\label{sec:io}
% ============================================================

\subsection{Parameter Parsing from \texttt{input.dat}}

All simulation parameters are read from a plain-text file \texttt{input.dat}
through the \texttt{io} module. The parser applies default values for any
missing keyword, strips inline comments (lines beginning with \texttt{\#}),
and performs basic range validation. The following parameters are controlled
from the input file:

\begin{table}[ht]
\centering
\caption{Simulation parameters read from \texttt{input.dat}.}
\label{tab:input_params}
\begin{tabular}{lll}
\toprule
Parameter & Default & Description \\
\midrule
\texttt{n\_carbons}   & 10       & Number of backbone carbon atoms \\
\texttt{n\_steps}     & 1{,}000{,}000 & Number of MC production steps \\
\texttt{temperature}  & 300.0    & Simulation temperature (K) \\
\texttt{conf\_type}   & 1        & Initial configuration mode (see Table~\ref{tab:conf_types}) \\
\texttt{explicit\_h}  & .false.  & Enable all-atom (explicit hydrogen) mode \\
\texttt{rng\_seed}    & 42       & Random number generator seed \\
\texttt{print\_interval} & 1{,}000 & Steps between output writes \\
\texttt{max\_delta}   & 0.5      & Maximum dihedral rotation per MC step (rad) \\
\midrule
\end{tabular}
\end{table}

Output is written in XYZ format: trajectories to
\texttt{traj\_\textit{label}.xyz}, energies to
\texttt{energy\_\textit{label}.dat}, and observables to
\texttt{observables\_\textit{label}.dat}, where the label encodes the key
simulation parameters (\texttt{n\_carbons}, \texttt{conf\_type},
\texttt{n\_steps}, \texttt{temperature}) for unambiguous identification.

% ============================================================